\chapter{Cztery ważne całki nieelementarne}

Wyniki dotyczą funkcji pierwotnych na niepustych otwartych przedziałach oraz ich lokalnych zespolonych odpowiedników. Wykluczenie funkcji pierwotnej nawet w zespolonym rozszerzeniu elementarnym wyklucza rzeczywisty wzór elementarny. Dla skończonej liczby jego podwyrażeń wybieramy wspólny zwykły punkt, unikając biegunów, rozgałęzień i zer argumentów logarytmów. Tam wszystkie wybrane gałęzie mają holomorficzne odpowiedniki, więc rzeczywisty wzór dawałby także wzór zespolony.

\section{Lokalny test biegunów}

Jeżeli funkcja wymierna $r$ ma w $a$ biegun rzędu $m\ge1$, to jej rozwinięcie zaczyna się od $c(x-a)^{-m}$, $c\ne0$. Pochodna zaczyna się od $-mc(x-a)^{-m-1}$ i ma biegun rzędu $m+1$. Funkcja wymierna bez biegunów skończonych jest wielomianem: mianownik nieskracalnego ułamka nie może mieć zera nad $\C$.

\begin{lemat}
Jeśli $h\in\C(x)$ jest regularne w $a$, a $s\in\C(x)$ ma tam biegun prosty, to $r'+hr=s$ nie ma rozwiązania $r\in\C(x)$.
\end{lemat}
\begin{proof}
Gdy $r$ jest regularne w $a$, lewa strona też jest regularna. Gdy $r$ ma biegun rzędu $m\ge1$, pochodna ma biegun rzędu $m+1$, a $hr$ co najwyżej $m$. Biegun rzędu $m+1\ge2$ nie może się skasować. W żadnym przypadku nie otrzymamy bieguna prostego.
\end{proof}

\section{Całka Gaussa: $\int e^{-x^2}\dd x$}

\begin{twierdzenie}
Funkcja $e^{-x^2}$ nie ma elementarnej funkcji pierwotnej nad $\C(x,e^{-x^2})$.
\end{twierdzenie}
\begin{proof}
Lemat 8.8 sprawdza przestępność $t=e^{-x^2}$ i brak nowych stałych. Mamy $Dt/t=-2x$. Twierdzenie 8.10 z $s=1$ wymagałoby $r\in\C(x)$ z
\[
r'-2xr=1.
\]
Każdy skończony biegun $r$ rzędu $m$ dawałby w $r'$ biegun rzędu $m+1$, którego $-2xr$ nie może usunąć: czynnik $-2x$ nie ma biegunów. Prawa strona jest regularna, więc $r$ nie ma żadnych biegunów skończonych i jest wielomianem.

Jeśli $r\ne0$ ma stopień $d\ge0$, to $-2xr$ ma stopień $d+1$, a $r'$ stopień mniejszy; dla stałego $r$ pochodna jest zerowa. Suma nie może być stałą $1$. Także $r=0$ nie jest rozwiązaniem. Kryterium daje brak elementarnej funkcji pierwotnej.
\end{proof}

Funkcja pierwotna istnieje jako funkcja określona całką; definiuje się na przykład funkcję
\[
\operatorname{erf}(x)=\frac{2}{\sqrt\pi}\int_0^x e^{-t^2}\dd t.
\]
Udowodniliśmy ograniczenie jej możliwej postaci, nie brak istnienia funkcji.

\subsection{Pełne obliczenie całki oznaczonej}

Niech $I=\int_0^\infty e^{-x^2}\dd x$. Całka jest zbieżna: na $[0,1]$ funkcja jest ograniczona, a dla $x\ge1$ mamy $e^{-x^2}\le e^{-x}$. Dodatniość funkcji, twierdzenie Tonellego i zamiana zmiennych na biegunowe dają
\[
I^2=\int_0^\infty\int_0^\infty e^{-(x^2+y^2)}\dd x\dd y
=\int_0^{\pi/2}\int_0^\infty e^{-r^2}r\dd r\dd\theta
=\frac\pi4.
\]
Ponieważ $I>0$, otrzymujemy $I=\sqrt\pi/2$. Argument dwuwymiarowy oblicza jedną liczbę, nie daje elementarnego wzoru na funkcję $x\mapsto\int_0^x e^{-t^2}\dd t$.

\section{Całka sinusowa: $\int \frac{\sin x}{x}\dd x$}

\begin{twierdzenie}
Funkcja $\sin x/x$, przedłużona wartością $1$ w zerze, nie ma elementarnej funkcji pierwotnej nad $\C(x,e^{ix})$.
\end{twierdzenie}
\begin{proof}
Niech $z=e^{ix}$. Przestępność i ciało stałych sprawdza lemat 8.8, zastosowany do $p(x)=ix$. Ponieważ $Dz=iz$, mamy
\[
\frac{\sin x}{x}=\frac{z-z^{-1}}{2ix}.
\]
Przy elementarnej funkcji pierwotnej dodatnia potęga $z$ wymuszałaby na mocy wniosku 8.11 rozwiązanie $r_1'+ir_1=1/(2ix)$ z $r_1\in\C(x)$. Dla $r=2ir_1$ dostalibyśmy
\[
r'+ir=\frac1x.
\]
To niemożliwe na mocy lematu 9.1 w punkcie $0$, z $h=i$. Nie trzeba już analizować potęgi $z^{-1}$.
\end{proof}

Nie twierdzimy, że $e^{ix}$ i $e^{-ix}$ są algebraicznie niezależne: ich iloczyn wynosi $1$. Użyliśmy dwóch potęg jednego przestępnego generatora. Brak rzeczywistej osobliwości samej funkcji $\sin x/x$ w zerze nie przeczy argumentowi: test biegunów dotyczy wymiernego równania wymuszonego przez współczynnik przy $z$, a nie całej funkcji. Oznaczamy
\[
\operatorname{Si}(x)=\int_0^x\frac{\sin t}{t}\dd t.
\]

\subsection{Dlaczego całka niewłaściwa wynosi połowę liczby $\pi$?}

Dla $\varepsilon>0$ wprowadźmy tłumienie
\[
I(\varepsilon)=\int_0^\infty e^{-\varepsilon x}\frac{\sin x}{x}\dd x.
\]
Całka jest bezwzględnie zbieżna. Na każdym zwartym przedziale dodatnich $\varepsilon$ pochodną funkcji podcałkowej dominuje funkcja $e^{-\delta x}$ dla pewnego $\delta>0$, więc wolno różniczkować pod znakiem całki. Otrzymujemy
\[
I'(\varepsilon)=-\int_0^\infty e^{-\varepsilon x}\sin x\dd x
=-\frac1{1+\varepsilon^2}.
\]
Ostatnią całkę obliczamy jako część urojoną $\int_0^\infty e^{-(\varepsilon-i)x}\dd x=1/(\varepsilon-i)$. Ponadto $|\sin x|/x\le1$ dla $x>0$ daje $|I(\varepsilon)|\le1/\varepsilon$, więc
\[
I(\varepsilon)=\int_\varepsilon^\infty\frac{\dd t}{1+t^2}
=\frac\pi2-\arctan\varepsilon,
\qquad\lim_{\varepsilon\downarrow0}I(\varepsilon)=\frac\pi2.
\]

Trzeba jeszcze wykazać, że ta granica jest całką niewłaściwą bez tłumienia. Dla $g_\varepsilon(x)=e^{-\varepsilon x}/x$, $\varepsilon\ge0$, funkcja $g_\varepsilon$ maleje do zera. Całkowanie przez części na $[A,B]$, $0<A<B$, daje
\[
\left|\int_A^B g_\varepsilon(x)\sin x\dd x\right|
\le g_\varepsilon(A)+g_\varepsilon(B)
+\int_A^B|g_\varepsilon'(x)|\dd x
=2g_\varepsilon(A)\le\frac2A.
\]
Dla $\varepsilon=0$ jest to kryterium Cauchy'ego zbieżności ogona. Oszacowanie jest jednostajne względem $\varepsilon\ge0$. Na $[0,A]$ funkcje podcałkowe zbiegają jednostajnie do $\sin x/x$, z wartością $1$ w zerze. Możemy więc przejść do granicy całek na tym odcinku, a ogony są dowolnie małe po wyborze dużego $A$. Stąd
\[
\int_0^\infty\frac{\sin x}{x}\dd x=\frac\pi2.
\]

\section{Logarytm całkowy: $\int \frac{\dd x}{\log x}$}

\begin{twierdzenie}
Na żadnym niepustym otwartym przedziale zawartym w $(0,1)$ lub $(1,\infty)$ funkcja $1/\log x$ nie ma elementarnej funkcji pierwotnej.
\end{twierdzenie}
\begin{proof}
Załóżmy, że istnieje elementarna funkcja pierwotna $F(x)$. Podstawmy $x=e^u$, wybierając lokalnie $u=\log x$. Reguła łańcuchowa daje
\[
\frac{\dd}{\dd u}F(e^u)=\frac{e^u}{u}.
\]
Podstawienie do skończonego wyrażenia elementarnego zachowuje jego elementarność: działania arytmetyczne, logarytmy i funkcje wykładnicze nadal są tego samego typu, a równania algebraiczne mają podstawione współczynniki. Dostawalibyśmy więc elementarną funkcję pierwotną $e^u/u$ nad $\C(u,e^u)$.

Twierdzenie 8.10, z $h=1$ i $s=1/u$, wymaga rozwiązania
\[
r'+r=\frac1u,\qquad r\in\C(u).
\]
Lemat 9.1 wyklucza je w punkcie $u=0$. Równanie dotyczy tożsamości funkcji wymiernych, dlatego jego biegun w zerze jest przeszkodą również wtedy, gdy wyjściowy przedział nie zawiera zera. Otrzymujemy sprzeczność.
\end{proof}

Do nazwania funkcji pierwotnej na $x>1$ wystarcza definicja
\[
\operatorname{Li}(x)=\int_2^x\frac{\dd t}{\log t}.
\]
Tutaj droga całkowania nie przechodzi przez $1$, więc nie potrzeba wartości głównej. Na $(0,1)$ można wybrać dowolny punkt początkowy w tym przedziale. Wybór stałej całkowania nie zmienia nieelementarności.

\section{Całka eliptyczna: $\int \frac{\dd x}{\sqrt{1-x^4}}$}

Ten przykład jest innego typu: całkowana funkcja jest algebraiczna, a nie wykładnicza. Zastosujemy samo twierdzenie Liouville'a oraz lokalne rzędy zer i biegunów. Nie potrzebujemy teorii rodzaju krzywej ani osobnego twierdzenia o całkach abelowych.

\subsection{Ciało funkcji i punkty nad nieskończonością}

Pracujemy w ciele $K=\C(x,y)$ z relacją $y^2=1-x^4$. Wielomian $1-x^4$ nie jest kwadratem funkcji wymiernej, bo ma proste zera, a kwadrat ma parzyste rzędy zer i biegunów. Dlatego $[K:\C(x)]=2$ i każdy element ma jednoznaczną postać $A(x)+B(x)y$, gdzie $A,B\in\C(x)$.

Derywację określają
\[
Dx=1,\qquad Dy=-\frac{2x^3}{y}.
\]
Jej stałe są równe $\C$, na mocy lematu 8.3. Funkcje z $K$ rozpatrujemy na krzywej $y^2=1-x^4$ wraz z punktami nad nieskończonością. Potrzebne lokalne współrzędne opisujemy jawnie.
\begin{itemize}
\item Gdy $y\ne0$, lokalną współrzędną jest $x-a$, gdzie $a$ jest współrzędną rozważanego punktu. Druga współrzędna jest lokalną gałęzią pierwiastka.
\item Przy $y=0$ mamy $x=a$ z $a^4=1$. Ponieważ pochodna $4x^3$ nie znika przy $x=a$, z twierdzenia o funkcji uwikłanej $x$ jest holomorficzną funkcją $y$ w otoczeniu tego punktu. Lokalną współrzędną jest $y$, a $x=a(1-y^2)^{1/4}$. Wybieramy gałąź pierwiastka o wartości $1$ w zerze, tak aby $x(0)=a$. W szczególności $x-a$ ma rząd $2$, a $y$ rząd $1$.
\item Przy nieskończoności wprowadzamy $s=1/x$, $w=y/x^2$, $w^2=s^4-1$. Dodajemy dwa punkty $(s,w)=(0,i)$ i $(0,-i)$, nie jeden. Ponieważ $2w\ne0$ w obu tych punktach, równanie $w^2=s^4-1$ wyznacza tam holomorficznie $w$ jako funkcję $s$. W obu lokalną współrzędną jest $s$, a $y=w/s^2$ ma biegun rzędu $2$. Te dwa punkty domykają opis nad nieskończonością: poza nimi i punktami o skończonym $x$ nie ma innych punktów wymaganych do badania zer i biegunów funkcji z $K$.
\end{itemize}
Zmiana między opisami jest dana wypisanymi funkcjami wymiernymi. W tych współrzędnych każda niezerowa funkcja z $K$ ma rozwinięcie Laurenta, więc można mówić o jej rzędzie $\ord_P(u)$ w punkcie $P$. Rząd jest dodatni dla zera, ujemny dla bieguna i zerowy dla funkcji regularnej i niezerowej w punkcie.

\begin{lemat}
Funkcja z $K$ bez biegunów we wszystkich wyżej opisanych punktach jest stała.
\end{lemat}
\begin{proof}
Niech $v=A+By$ nie ma biegunów. Inwolucja $\tau(x,y)=(x,-y)$ permutuje punkty krzywej, także dwa punkty nad nieskończonością. Dlatego $\tau(v)=A-By$ również nie ma biegunów. Stąd ani $A=(v+\tau(v))/2$, ani $By=(v-\tau(v))/2$ nie mają biegunów.

Funkcja wymierna $A(x)$ nie może mieć bieguna w żadnym skończonym $x=a$, bo dawałoby to biegun w punkcie nad $a$. Jest więc wielomianem. Nie może też mieć dodatniego stopnia, bo wtedy ma biegun przy $s=0$. Zatem $A$ jest stała.

Jeśli $B$ miałoby biegun rzędu $m\ge1$ w punkcie $a$ z $1-a^4\ne0$, to $By$ miałoby tam biegun, ponieważ $y$ jest regularne i niezerowe. Dla $a^4=1$ rząd $By$ wynosiłby $1-2m<0$, co też daje biegun. Funkcja $B$ nie ma więc skończonych biegunów i jest wielomianem. Jeśli $B\ne0$ ma stopień $d$, to przy nieskończoności $By$ ma biegun rzędu $d+2$. Stąd $B=0$, a $v=A$ jest stała.
\end{proof}

\subsection{Forma różniczkowa bez biegunów}

Forma różniczkowa ma w lokalnej współrzędnej $z$ postać $a(z)\dd z$ i zmienia się według reguły łańcuchowej. Jest holomorficzna, jeśli $a(z)$ nie ma bieguna. Dla $v\in K$ mamy $\dd v=(Dv)\dd x$. Sprawdźmy formę
\[
\omega=\frac{\dd x}{y}.
\]
Przy $y\ne0$ jest regularna i niezerowa. Przy punkcie rozgałęzienia różniczkowanie $y^2=1-x^4$ daje $2y\dd y=-4x^3\dd x$, więc $\omega=-\dd y/(2x^3)$. Ponieważ tam $x^4=1$, współczynnik jest regularny i niezerowy. Przy nieskończoności
\[
\omega=\frac{-s^{-2}\dd s}{w/s^2}=-\frac{\dd s}{w},
\]
co również jest regularne i niezerowe przy $w=\pm i$. Forma $\omega$ nie ma więc bieguna w żadnym punkcie i nie jest zerowa.

\subsection{Dlaczego nie może to być pochodna elementarna?}

\begin{twierdzenie}
Element $1/y\in K$, $y^2=1-x^4$, nie ma elementarnej funkcji pierwotnej. W szczególności $1/\sqrt{1-x^4}$ nie ma takiej funkcji pierwotnej na $(-1,1)$.
\end{twierdzenie}
\begin{proof}
Gdyby funkcja pierwotna była elementarna, twierdzenie Liouville'a zastosowane do $K$ i pomnożone przez $\dd x$ dawałoby
\begin{equation}\label{eq:elliptic-liouville}
\omega=\dd v+\sum_{j=1}^r c_j\frac{\dd u_j}{u_j},\qquad
v\in K,\quad u_j\in K^\times,\quad c_j\in\C.
\end{equation}
Wybrana zwykła gałąź daje zanurzenie różniczkowe $K$ w ciele zarodków meromorficznych. Jest ono różnowartościowe, jak każdy homomorfizm ciał. Dlatego równość uzyskana z lokalnej funkcji pierwotnej jest tożsamością w $K$; możemy badać ją we wszystkich punktach. Nie zakładamy globalnego przedłużenia samej hipotetycznej funkcji pierwotnej.

\emph{Część dokładna.} Gdyby $v$ miało biegun rzędu $m\ge1$ w lokalnej współrzędnej $z$, to $\dd v$ miałoby biegun rzędu $m+1\ge2$. Tymczasem dla $u=z^k h(z)$, z $h$ regularną i $h(0)\ne0$, mamy
\[
\frac{\dd u}{u}=k\frac{\dd z}{z}+\frac{\dd h}{h}.
\]
Pochodne logarytmiczne mają co najwyżej bieguny proste. Nie mogą usunąć bieguna $\dd v$ rzędu co najmniej $2$, a $\omega$ nie ma biegunów. Zatem $v$ nie ma biegunów i na mocy lematu 9.5 jest stałe. Część $\dd v$ znika.

\emph{Część logarytmiczna.} Współczynnik przy $\dd z/z$, zwany residuum, w $\dd u/u$ jest równy całkowitej liczbie $\ord_P(u)$. Ponieważ $\omega$ jest holomorficzna, dla każdego punktu $P$ mamy
\[
\sum_j c_j\ord_P(u_j)=0.
\]
Wybierzmy bazę $\gamma_1,\dots,\gamma_s$ nad $\Q$ przestrzeni rozpiętej przez skończenie wiele $c_j$ i zapiszmy $c_j=\sum_i q_{ij}\gamma_i$, $q_{ij}\in\Q$. Niezależność nad $\Q$ daje
\[
\sum_j q_{ij}\ord_P(u_j)=0\qquad\text{dla każdego }i\text{ i }P.
\]
Wybierzmy wspólny dodatni mianownik $N$ wszystkich $q_{ij}$ i połóżmy $U_i=\prod_j u_j^{Nq_{ij}}$. Wykładniki są całkowite, więc $U_i$ jest zwykłym niezerowym elementem $K$, a nie wielowartościową potęgą. We wszystkich punktach ma rząd zero. Nie ma biegunów, więc jest stały na mocy lematu 9.5. Stąd
\[
\sum_j c_j\frac{\dd u_j}{u_j}
=\frac1N\sum_i\gamma_i\frac{\dd U_i}{U_i}=0.
\]
Jeśli wszystkie $c_j$ są zerowe, wniosek jest natychmiastowy bez wyboru bazy. Tożsamość \eqref{eq:elliptic-liouville} dawałaby więc $\omega=0$, wbrew bezpośredniemu sprawdzeniu.

Istotne były całkowite rzędy zer i biegunów funkcji $u_j$. Samo znikanie residuów nie dowodzi, że dowolna forma jest zerowa: właśnie $\omega$ jest niezerową formą bez residuów. Udowodniliśmy znikanie konkretnej kombinacji pochodnych logarytmicznych, wykorzystując bazę nad $\Q$.
\end{proof}

\section{Przykłady dla kontrastu}

Niewielka zmiana całkowanej funkcji może przywrócić elementarność:
\[
\int x e^{x^2}\dd x=\frac12e^{x^2}+C,\qquad
\int\frac{\dd x}{x\log x}=\log|\log x|+C,
\qquad
\int\frac{2x}{x^2+1}\dd x=\log(x^2+1)+C.
\]
Drugi wzór jest rzeczywisty na każdym z przedziałów $(0,1)$ i $(1,\infty)$; w wersji zespolonej wybieramy lokalną gałąź $\log(\log x)$. Wszystkie równości sprawdzamy przez różniczkowanie. W pierwszej równanie $r'+2xr=x$ ma rozwiązanie $r=1/2$; w pozostałych występują pochodne logarytmiczne.

\begin{twierdzenie}
Każda funkcja wymierna nad $\C$ ma lokalną elementarną funkcję pierwotną.
\end{twierdzenie}
\begin{proof}
Odejmując skończone części główne rozwinięć przy wszystkich biegunach, otrzymujemy funkcję wymierną bez biegunów skończonych, czyli wielomian. Daje to rozkład
\[
f(x)=P(x)+\sum_a\sum_{k=1}^{m_a}\frac{c_{a,k}}{(x-a)^k}.
\]
Wielomian całkujemy wyraz po wyrazie. Dla $k\ge2$ mamy
\[
\int\frac{c_{a,k}}{(x-a)^k}\dd x
=-\frac{c_{a,k}}{(k-1)(x-a)^{k-1}},
\]
a dla $k=1$ funkcją pierwotną jest $c_{a,1}\log(x-a)$ z lokalną gałęzią. Skończona suma tych wyrażeń jest elementarna i ma pochodną $f$.
\end{proof}


